\documentclass[10pt, a4paper]{article}
\usepackage[a4paper, top=0.85cm, bottom=0.85cm, left=1.0cm, right=1.0cm]{geometry}
\ifx\Umathchar\undefined
  \usepackage[utf8]{inputenc}
  \usepackage[T1]{fontenc}
  \usepackage[spanish, es-nodecimaldot]{babel}
  \usepackage{graphicx}
\else
  \usepackage{fontspec}
  \usepackage[spanish, english, bidi=basic, provide=*]{babel}
  \babelprovide[import, onchar=ids fonts]{spanish}
  \babelprovide[import, onchar=ids fonts]{english}
  \babelfont{rm}{Noto Sans}
\fi
\usepackage{indentfirst}
\usepackage{enumitem}
\setlist[itemize]{label=-, nosep, leftmargin=*}
\usepackage{amsmath}
\usepackage{booktabs}
\usepackage{float}
\usepackage{caption}
\captionsetup{font=scriptsize, labelfont=bf, skip=2pt}

\newcommand{\mostrarfigura}[4]{%
  \IfFileExists{#1.#2}{%
    \csname includegraphics\endcsname[width=\linewidth,height=#3,keepaspectratio]{#1.#2}%
  }{%
    \framebox{\parbox{0.93\linewidth}{\centering\vspace{0.18cm}\scriptsize\textbf{[#4]}\vspace{0.18cm}}}%
  }%
}

\setlength{\parindent}{1.8em}
\setlength{\parskip}{1.5pt}
\renewcommand{\baselinestretch}{0.94}

\makeatletter
\renewcommand\section{\@startsection {section}{1}{\z@}%
  {-1.0ex \@plus -0.2ex \@minus -0.1ex}%
  {0.35ex \@plus 0.1ex}%
  {\normalfont\normalsize\bfseries}}
\makeatother

\title{\vspace{-1.5cm}\textbf{\large Informe de Avance 2 -- Concentración Mensual de $\text{MP}_{2,5}$ en La Florida (2002--2026)}\vspace{-0.35cm}}
\author{\footnotesize \textbf{Autores:} Luna Garcés, Tomás Pérez, Pilar Vallotton y Julián Vargas \quad|\quad \textbf{Curso:} EYP290I-1 Series de Tiempo\vspace{-0.35cm}}
\date{\scriptsize 16 de noviembre de 2026}

\begin{document}
\maketitle
\thispagestyle{empty}
\pagestyle{empty}
\footnotesize

En el presente informe se analiza la concentración de material particulado fino ($\text{MP}_{2,5}$, en $\mu\text{g/m}^3$) medida por la estación La Florida del Sistema de Información Nacional de Calidad del Aire (SINCA, Ministerio del Medio Ambiente). La base original contiene registros diarios desde enero de 2002 hasta agosto de 2026. Para el análisis se agregaron los registros validados en promedios mensuales, obteniendo una serie regular de $n = 296$ observaciones. El objetivo es determinar si existen patrones de tendencia y estacionalidad en la serie (es decir, si el nivel de contaminación cambia estructuralmente con los años y si se repite un ciclo anual), y evaluar qué tan bien pueden predecirse sus valores futuros mediante modelos de series de tiempo.

\section*{1. Datos, Partición Temporal y Limpieza}

A los registros diarios de $\text{MP}_{2,5}$ se les calcularon promedios mensuales utilizando únicamente los valores \textbf{validados} por el SINCA (\texttt{na.action = na.pass} para preservar el índice temporal), excluyendo los registros preliminares, los no validados y la columna de hora (compuesta únicamente por ceros). De los 9.000 días disponibles en el período de estudio, 8.304 ($92,3$\%) corresponden a datos validados, conformando una serie mensual de $n = 296$ observaciones entre enero de 2002 y agosto de 2026 (mes que cuenta con 22 días medidos). Dentro de la serie, 14 meses intermedios carecen de registros validados: enero a septiembre de 2013 (9 meses) y enero a mayo de 2017 (5 meses).

Con el fin de evaluar la capacidad predictiva fuera de muestra sin filtrar información futura (\textit{data leakage}), la serie se dividió \textbf{antes} de imputar en un conjunto de entrenamiento con los primeros $n_{\text{train}} = 266$ meses ($90$\%, enero de 2002 a febrero de 2024, índice decimal $2024,083$) y un conjunto de validación con los $n_{\text{val}} = 30$ meses más recientes ($10$\%, marzo de 2024 a agosto de 2026, índice $2024,167$, sin datos faltantes). Los 14 meses faltantes del entrenamiento se imputaron mediante el suavizador de Kalman (\texttt{KalmanSmooth}) sobre un modelo de espacio de estados $\text{SARIMA}(1,0,0)(0,1,1)_{12}$ ajustado a la serie logarítmica centrada, proyectando cada observación faltante sobre el pasado y futuro dentro del entrenamiento ($E[\log X_s \mid \log X_{1:n_{\text{train}}}]$). Trabajar en escala logarítmica garantiza que las imputaciones respeten el soporte positivo al aplicar $\exp()$ y reproduce fielmente la fase del ciclo anual medida por los registros preliminares del sensor (Figura 1).

\begin{figure}[H]
    \centering
    \begin{minipage}[t]{0.49\linewidth}
        \centering
        \mostrarfigura{Train_validation}{png}{2.65cm}{Figura 1: Serie mensual, corte 90/10 e imputación de Kalman}
        \caption{Serie mensual de $\text{MP}_{2,5}$, partición 90/10 e imputación de Kalman.}
        \label{fig:imputacion}
    \end{minipage}\hfill
    \begin{minipage}[t]{0.49\linewidth}
        \centering
        \mostrarfigura{ACF_PACF_Periodograma}{png}{2.65cm}{Figura 2: ACF, PACF y periodograma suavizado}
        \caption{ACF, PACF y periodograma suavizado de $\log(\text{MP}_{2,5})$.}
        \label{fig:exploratorio}
    \end{minipage}
\end{figure}
\vspace{-0.25cm}

\section*{2. Análisis Exploratorio, Transformación y Estacionariedad}

Al observar la evolución histórica de la serie mensual de $\text{MP}_{2,5}$, se distingue un comportamiento estacional pronunciado, con alzas recurrentes de contaminación en otoño e invierno ($45\text{--}62\ \mu\text{g/m}^3$) asociadas a mayores emisiones por calefacción y episodios de mala ventilación en la cuenca de Santiago, seguidas de caídas en verano ($10\text{--}18\ \mu\text{g/m}^3$), además de una disminución paulatina en los valles de verano a lo largo de las dos décadas. En la Figura 2, la ACF de $\log(X_t)$ exhibe una oscilación sinusoidal persistente con picos significativos en múltiplos anuales ($12, 24, 36$ y $48$), mientras que la PACF muestra correlaciones parciales activas en los primeros rezagos y en torno al rezago 12. Tanto el periodograma crudo como su estimación consistente suavizada con splines (\texttt{smooth.periodogram}, $\text{spar} = 0,5$) concentran la energía espectral en la frecuencia $\omega = 0,5197\text{ rad}$ ($j = 22$ en la grilla $2\pi j/266$), equivalente a un período exacto de $2\pi/0,5197 = 12,09 \approx 12$ meses.

\begin{figure}[H]
    \centering
    \begin{minipage}[t]{0.49\linewidth}
        \centering
        \mostrarfigura{Normal_loglik_escalalog}{png}{2.65cm}{Figura 3: Comparación del CV anual (2002--2023)}
        \caption{Comparación del CV anual (2002--2023) en las tres escalas.}
        \label{fig:transformacion}
    \end{minipage}\hfill
    \begin{minipage}[t]{0.49\linewidth}
        \centering
        \mostrarfigura{Transf_Diferenciacion}{png}{2.65cm}{Figura 4: Serie, ACF y periodograma tras diferenciar}
        \caption{Serie, ACF y periodograma tras aplicar la diferencia estacional.}
        \label{fig:diferenciacion}
    \end{minipage}
\end{figure}
\vspace{-0.25cm}

Dado que la amplitud estacional escala con el nivel anual de contaminación (heterocedasticidad multiplicativa), se evaluó la transformación de Box-Cox sobre el entrenamiento, obteniendo $\hat{\lambda} = 0,55$ mediante el criterio de Guerrero y $\hat{\lambda} = 0,20$ por perfil de log-verosimilitud (\texttt{loglik}). Al comparar el coeficiente de variación ($\text{CV} = s/\bar{x}$) en los años calendario completos del entrenamiento (2002--2023) en la Figura 3, se comprueba que la escala original presenta alta inestabilidad ($\text{CV}$ entre $0,26$ y $0,56$), mientras que la transformación logarítmica $Y_t = \log(X_t)$ ($\lambda = 0$) logra los menores niveles de CV ($0,08\text{--}0,17$) y permite interpretar la deriva como variación porcentual. Siguiendo la jerarquía de diferenciación estacional vista en clases, la función \texttt{nsdiffs} sobre $Y_t$ indicó aplicar una diferencia estacional ($D = 1$, período 12), definiendo $W_t = (1 - B^{12})\log(X_t)$. Como el operador $(1 - B^{12})$ elimina el ciclo anual y transforma cualquier tendencia lineal en una constante (\texttt{drift}), las pruebas \texttt{ndiffs} (criterios ADF, KPSS y PP) evaluadas sobre $W_t$ coincidieron en que no se requiere diferenciación regular ($d = 0$), lo cual fue corroborado por el test de Dickey-Fuller aumentado ($\text{ADF} = -5,1671$, rezago 6, $p < 0,01$). En la Figura 4 se confirma que tras aplicar $(1 - B^{12})$ desaparece el máximo estacional del periodograma y la ACF queda acotada a memoria corta.

\section*{3. Modelo Preliminar, Depuración y Diagnóstico de Residuos}

Con $d = 0$ y $D = 1$, se ajustaron por máxima verosimilitud cinco modelos competidores sobre $X_{\text{train}}$ ($\lambda = 0$): (1) el modelo Airline $\text{SARIMA}(0,1,1)(0,1,1)_{12}$; (2) un modelo con doble diferencia $\text{SARIMA}(1,1,1)(0,1,1)_{12}$ para verificar el costo de sobrediferenciar ($d = 1$); (3) un $\text{SARIMA}(1,0,1)(0,1,1)_{12}$ con \texttt{drift}; (4) el modelo parsimonioso $\text{SARIMA}(1,0,0)(0,1,1)_{12}$ con \texttt{drift} tras descartar el coeficiente $\text{MA}(1)$ no significativo ($p = 0,43$); y (5) el modelo depurado $\text{SARIMA}(1,0,[15])(0,1,1)_{12}$ con \texttt{drift}, que utiliza el argumento \texttt{fixed} para fijar $\theta_1 = \dots = \theta_{14} = 0$ y estimar $\theta_{15}$ con el fin de limpiar la correlación residual observada en el rezago 15.

\begin{table}[H]
\centering
\scriptsize
\setlength{\tabcolsep}{2.4pt}
\caption{Selección de modelos en entrenamiento ($n_{\text{train}}=266$) y métricas de validación fuera de muestra ($n_{\text{val}}=30$, mar-2024 a ago-2026).}
\begin{tabular}{lcccccccccc}
\toprule
\textbf{Modelo Competidor} & \textbf{AIC} & \textbf{AICc} & \textbf{BIC} & \textbf{$p$-BL(12)} & \textbf{$p$-BL(15)} & \textbf{$p$-BL(24)} & \textbf{RMSE} & \textbf{MAE} & \textbf{MAPE (\%)} & \textbf{Cob. $95$\%} \\
\midrule
Airline $(0,1,1)(0,1,1)_{12}$ & $-201,30$ & $-201,21$ & $-190,70$ & $0,000$ & $0,000$ & $0,000$ & $4,84$ & $3,83$ & $17,58$ & $100,0$\% \\
Sobredif. $(1,1,1)(0,1,1)_{12}$ & $-220,16$ & $-219,99$ & $-206,02$ & $0,416$ & $0,018$ & $0,063$ & $4,80$ & $3,79$ & $16,97$ & $96,67$\% \\
$\text{SARIMA}(1,0,1)(0,1,1)_{12}$ + drift & $-226,19$ & $-225,95$ & $-208,50$ & $0,416$ & $0,009$ & $0,043$ & $4,93$ & $3,54$ & $14,18$ & $96,67$\% \\
$\text{SARIMA}(1,0,0)(0,1,1)_{12}$ + drift & $-226,36$ & $-226,20$ & $-212,21$ & $0,230$ & $0,005$ & $0,033$ & $4,95$ & $3,53$ & $14,09$ & $96,67$\% \\
\textbf{Depurado $(1,0,[15])(0,1,1)_{12}$ + drift} & $\mathbf{-231,22}$ & $\mathbf{-230,98}$ & $\mathbf{-213,53}$ & $\mathbf{0,170}$ & $\mathbf{0,024}$ & $\mathbf{0,075}$ & $\mathbf{4,85}$ & $\mathbf{3,36}$ & $\mathbf{13,28}$ & $\mathbf{96,67}$\textbf{\%} \\
\bottomrule
\end{tabular}
\end{table}
\vspace{-0.2cm}

Como muestra la Tabla 1, imponer $d = 1$ deteriora el AICc y el BIC en más de 6 unidades por sobrediferenciación frente a mantener $d = 0$ con \texttt{drift}. Al comparar los dos mejores candidatos en la Tabla 2, todos sus parámetros resultan altamente significativos bajo normalidad asintótica del estimador de máxima verosimilitud ($\sqrt{n}(\hat{\theta}-\theta_0) \xrightarrow{d} \text{N}(0, \Gamma^{-1})$). La ecuación estimada del modelo depurado óptimo es $(1 - 0,4693 B)(1 - B^{12})(\log X_t + 0,00194 t) = (1 - 0,2079 B^{15})(1 - 0,8120 B^{12})\hat{\epsilon}_t$. En el contexto ambiental de La Florida, $\hat{\phi}_1 = 0,4693$ refleja una persistencia mensual del $46,9$\% ante anomalías atmosféricas; $\hat{\Theta}_1 = -0,8120$ captura la estabilización del perfil invernal entre años sucesivos; y la deriva $\widehat{\text{drift}} = -0,00194$ cuantifica una reducción estructural media de $0,194$\% mensual ($\approx 2,30$\% anual) en la concentración base de $\text{MP}_{2,5}$.

\begin{table}[H]
\centering
\scriptsize
\setlength{\tabcolsep}{2.1pt}
\caption{Inferencia asintótica (\texttt{coeftest}) de los modelos $\text{SARIMA}(1,0,0)(0,1,1)_{12}$ y $\text{SARIMA}(1,0,[15])(0,1,1)_{12}$ con \texttt{drift}, y outliers detectados ($|Z|>2,5$).}
\begin{tabular}{lccccclccc}
\toprule
& \multicolumn{2}{c}{\textbf{Parsimonioso $(1,0,0)(0,1,1)_{12}$}} & \multicolumn{3}{c}{\textbf{Depurado $(1,0,[15])(0,1,1)_{12}$}} & \multicolumn{4}{c}{\textbf{Outliers en Entrenamiento ($|Z| > 2,5$)}} \\
\cmidrule(lr){2-3} \cmidrule(lr){4-6} \cmidrule(lr){7-10}
\textbf{Parámetro} & \textbf{Estim. (EE)} & \textbf{Valor-$p$} & \textbf{Estim. (EE)} & \textbf{Estad. $z$} & \textbf{Valor-$p$} & \textbf{Fecha} & \textbf{$\text{MP}_{2,5}$} & \textbf{$Z_{\text{res}}$} & \textbf{Causa / Diagnóstico} \\
\midrule
$\phi_1$ (\texttt{ar1}) & $0,4458\ (0,0575)$ & $<0,0001$ & $0,4693\ (0,0570)$ & $8,233$ & $<0,0001$ & \texttt{2014-02} & $15,41$ & $-2,53$ & Valle estival limpio \\
$\theta_{15}$ (\texttt{ma15}) & Fijo en $0$ & --- & $-0,2079\ (0,0806)$ & $-2,578$ & $0,0099$ & \texttt{2022-05} & $50,39$ & $+2,98$ & Episodio crítico otoño \\
$\Theta_1$ (\texttt{sma1}) & $-0,8192\ (0,0485)$ & $<0,0001$ & $-0,8120\ (0,0510)$ & $-15,925$ & $<0,0001$ & \texttt{2023-02} & $27,54$ & $+3,89$ & Incendios centro-sur \\
\texttt{drift} & $-0,00188\ (0,00034)$ & $<0,0001$ & $-0,00194\ (0,00029)$ & $-6,634$ & $<0,0001$ & \multicolumn{4}{l}{\textit{Shapiro-Wilk (\texttt{res.clean}):} $p = 0,5444$} \\
\bottomrule
\end{tabular}
\end{table}
\vspace{-0.2cm}

El diagnóstico gráfico de residuos se presenta en los cuatro paneles de la Figura 5: (a) el primer panel verifica que la raíz inversa AR y las 12 raíces inversas SMA se sitúan estrictamente dentro del círculo unitario (proceso causal e invertible); (b) al contrastar el segundo panel ($\text{SARIMA}(1,0,0)(0,1,1)_{12}$, donde el rezago 15 rechaza blancura y arrastra $p_{\text{BL}(24)} = 0,033$) contra el tercer panel ($\text{SARIMA}(1,0,[15])(0,1,1)_{12}$), se observa que $\theta_{15}$ elimina el pico en el rezago 15 de la ACF residual y eleva el valor-$p$ de Box-Ljung en el rezago 24 a $0,075 > 0,05$; y (c) en el cuarto panel, tras aislar los tres episodios atípicos de la Tabla 2 (dominados por el humo de incendios forestales de febrero de 2023, $Z = +3,89$), los residuos limpios presentan simetría ($\text{asimetría} = 0,021$), curtosis normal ($\kappa = 2,621$) y ajuste sobre la recta del QQ-plot (Shapiro-Wilk $p = 0,544$). Asimismo, la ACF de $\hat{e}_t^2$ y el test de Box-Ljung sobre $\hat{e}_t^2$ ($p = 0,196$ en rezago 12 y $p = 0,497$ en rezago 24) descartan heterocedasticidad condicional ARCH.

\begin{figure}[H]
    \centering
    \begin{minipage}[t]{0.24\linewidth}
        \centering
        \mostrarfigura{AR_roots}{png}{2.25cm}{Figura 5a: Raíces AR y MA}
    \end{minipage}\hfill
    \begin{minipage}[t]{0.24\linewidth}
        \centering
        \mostrarfigura{Residuos1}{png}{2.25cm}{Figura 5b: ts.diag (1,0,0)}
    \end{minipage}\hfill
    \begin{minipage}[t]{0.24\linewidth}
        \centering
        \mostrarfigura{Residuos2}{png}{2.25cm}{Figura 5c: ts.diag ma15}
    \end{minipage}\hfill
    \begin{minipage}[t]{0.24\linewidth}
        \centering
        \mostrarfigura{Distribución_QQplot_residuos}{png}{2.25cm}{Figura 5d: QQ-plot y ARCH}
    \end{minipage}
    \caption{Diagnóstico integral de residuos: (a) raíces inversas, (b) \texttt{ts.diag} modelo $(1,0,0)$, (c) \texttt{ts.diag} modelo depurado \texttt{ma15} y (d) normalidad y ACF de $\hat{e}_t^2$.}
    \label{fig:diagnostico}
\end{figure}
\vspace{-0.25cm}

\section*{4. Validación Fuera de Muestra, Back-test Móvil y Próximos Pasos}

En la validación sobre los $h = 30$ meses de \texttt{X\_val} (marzo de 2024 a agosto de 2026, Figura 6), el modelo depurado (panel a) captura la trayectoria estacional con un $\text{MAPE} = 13,28$\% y una cobertura del $96,67$\% (29 de 30 meses dentro del intervalo al $95$\%, escapando solo el mínimo de febrero de 2026: observado $10,39\ \mu\text{g/m}^3$ frente al límite inferior $11,05\ \mu\text{g/m}^3$ del IC al $95$\%). Además, el panel (b) confirma la propiedad teórica vista en clases: con $d = 0$ y \texttt{drift} la amplitud de la banda converge, mientras que con $d = 1$ (sobrediferenciado) se abre excesivamente sin mejorar la cobertura. Como $\exp(\hat{\mu}_h)$ estima la mediana condicional, en el panel (c) se verifica que al aplicar la corrección por sesgo log-normal de la función generadora de momentos ($\hat{X}_{t+h} = \exp(\hat{\mu}_h + \hat{\sigma}_h^2/2)$, \texttt{biasadj = TRUE}) para obtener la media y al evaluar el predictor operativo a 1 paso ($h = 1$), los peaks invernales proyectados suben y el error baja de $\text{RMSE} = 4,85$ a $\mathbf{4,75\ \mu\text{g/m}^3}$ y $4,79\ \mu\text{g/m}^3$ respectivamente (Tabla 3). Finalmente, en el \textit{back-test} móvil a 12 pasos (panel d), el RMSE oscila entre $4,57$ ($h = 6$) y $5,21\ \mu\text{g/m}^3$ ($h = 12$) conforme crece la incertidumbre teórica $\sigma^2(h) = \sigma_\epsilon^2 \sum_{j=0}^{h-1}\psi_j^2$, manteniendo la cobertura empírica clavada en torno al nivel nominal del $95$\% ($93,1$\%--$96,4$\%). Como próximos pasos para el póster final, se incorporarán intervenciones de pulso para modelar eventos de incendios forestales (\texttt{2023-02}) y regresores meteorológicos (SARIMAX) que permitan anticipar los máximos invernales extremos.

\begin{table}[H]
\centering
\scriptsize
\setlength{\tabcolsep}{2.3pt}
\caption{Comparación de predictores en validación ($n_{\text{val}}=30$) y desempeño del \textit{back-test} móvil según el horizonte de pronóstico $h = 1, \dots, 12$ meses.}
\begin{tabular}{lcccclcccccc}
\toprule
\textbf{Estrategia en Validación ($h=30$)} & \textbf{RMSE} & \textbf{MAE} & \textbf{MAPE (\%)} & & \textbf{Métrica Back-test} & $\mathbf{h=1}$ & $\mathbf{h=2}$ & $\mathbf{h=4}$ & $\mathbf{h=6}$ & $\mathbf{h=9}$ & $\mathbf{h=12}$ \\
\midrule
Proyección 30m Mediana ($\exp(\hat{\mu}_h)$) & $4,85$ & $3,36$ & $13,28$ & & RMSE ($\mu\text{g/m}^3$) & $4,791$ & $4,961$ & $4,825$ & $4,573$ & $4,859$ & $5,213$ \\
Proyección 30m Media (\texttt{biasadj=TRUE}) & $\mathbf{4,75}$ & $3,37$ & $13,68$ & & MAPE (\%) & $13,61$ & $14,36$ & $13,89$ & $13,29$ & $14,55$ & $15,64$ \\
Pronóstico a 1 paso ($h=1$, filtro Kalman) & $4,79$ & $3,37$ & $13,61$ & & Cobertura IC $95$\% & $93,3$\% & $93,1$\% & $96,3$\% & $96,0$\% & $95,5$\% & $94,7$\% \\
\bottomrule
\end{tabular}
\end{table}
\vspace{-0.25cm}

\begin{figure}[H]
    \centering
    \begin{minipage}[t]{0.24\linewidth}
        \centering
        \mostrarfigura{Pronostico_Validación_SARIMA}{png}{2.25cm}{Figura 6a: Pronóstico y Zoom}
    \end{minipage}\hfill
    \begin{minipage}[t]{0.24\linewidth}
        \centering
        \mostrarfigura{Bandas_prediccion}{png}{2.25cm}{Figura 6b: Bandas d=0 vs d=1}
    \end{minipage}\hfill
    \begin{minipage}[t]{0.24\linewidth}
        \centering
        \mostrarfigura{Validación_30m_vs_Pronostico1paso}{png}{2.25cm}{Figura 6c: Media vs 1 paso}
    \end{minipage}\hfill
    \begin{minipage}[t]{0.24\linewidth}
        \centering
        \mostrarfigura{RMSE_calibración}{png}{2.25cm}{Figura 6d: RMSE y Cobertura}
    \end{minipage}
    \caption{Validación (2024--2026): (a) pronóstico y zoom, (b) bandas al $95$\% ($d=0$ vs $d=1$), (c) proyección media vs 1 paso y (d) \textit{back-test} móvil.}
    \label{fig:validacion}
\end{figure}
\vspace{-0.3cm}
\begin{flushright}
    \scriptsize \textit{Muchas gracias.}
\end{flushright}
\begin{flushright}
    \scriptsize \textit{\textbf{Declaración de uso de IA:} Fue utilizada para ajustar  la redacción y la estética del informe.}
\end{flushright}

\end{document}